\documentclass[10pt,a4paper]{amsart}
\usepackage[margin=19mm]{geometry}
\usepackage{amsmath,amssymb,mathtools}
\usepackage[scr=rsfs,cal=euler]{mathalfa}
\usepackage{xfrac}
\usepackage[unicode,hidelinks,bookmarks=false]{hyperref}
\newcommand{\ie}{i.e.\ }
\newcommand{\cf}{cf.\ }
\newcommand{\resp}{respectively\ }
\newcommand{\Subsection}{\subsection}
\providecommand{\nobreakdash}{\nobreak\hbox{-}}
\setlength{\emergencystretch}{2em}
\multlinegap0pt
\usepackage{mdframed}
\usepackage{mdframed}
\usepackage{mdframed}
\usepackage{mdframed}
\usepackage{mdframed}
\usepackage{csquotes}
\usepackage{amssymb}
\usepackage{dsfont}
\usepackage{xcolor}
\usepackage{braket}
\usepackage{mathtools}
\usepackage{cancel}
\usepackage{tikz}
\newtheorem{lem}{}
\usepackage{cleveref}
\crefname{lem}{}{s}
\newcommand{\RR}{\mathbb{R}}
\newcommand{\TT}{\mathbb{T}}
\newcommand{\dif}{\operatorname{d}\!} % shorten the space afterwards
\newcommand{\ee}{\operatorname{e}}
\newcommand{\eq}[1]{\begin{align*}#1\end{align*}}
\newcommand{\eql}[1]{\begin{align}#1\end{align}}
\newcommand{\ii}{\operatorname{i}}
\newcommand{\tr}{\operatorname{tr}}
\newcommand{\ve}{\varepsilon}
\title{Euler Topological Metals in 1D}
\author{Yichen Hu, Jacob Shapiro}
\date{Source: arXiv:2608.27554v1 (CC BY 4.0)}
\begin{document}
\maketitle

\noindent\emph{Source and licence.} Euler Topological Metals in 1D. Version arXiv:2608.27554v1 (CC BY 4.0), distributed by arXiv under the Creative Commons Attribution 4.0 International licence, http://creativecommons.org/licenses/by/4.0/, as recorded in the arXiv licence metadata for this exact version and shown on the abstract page https://arxiv.org/abs/2608.27554v1. Full licence notice: \texttt{licenses/arxiv-2608.27554-CC-BY-4.0.txt}.\par\smallskip

\noindent\textit{Editorial scope.} Counted unit: the article's lem formula (source0.tex lines 1255--1286). The article's complete original proof of it is delivered with it (source0.tex lines 1288--1391, one \texttt{proof} environment of 104 lines). No mathematics is written, repaired or re-derived: the statement and the proof are the article's own lines, in the article's own notation, and no retained line is substituted or re-flowed.  Retained uncounted as the source-local context the proof needs: lines 1215--1222.  Nothing was omitted from any retained span, so the package carries no omission record.  No retained line contains a citation, so the wrapper carries no bibliography and no bibliography entry.  No retained line contains a figure environment, an image-inclusion command or drawing code, so no image is delivered and the package declares no asset dependency.  Editorial formatting only, in addition: an AMS \texttt{amsart} preamble that restates the article's own declarations and macros used by the retained lines, namely the environments \texttt{lem} on separate counters, together with the standard packages those lines need.  The retained environments print in this document as Lemma 1 in order of appearance, whereas the article prints the same items with the numbering of its own full document.  The complete native source of the article as distributed in the arXiv e-print for this exact version is bundled unchanged as \texttt{sources/208722/source0.tex}.
\begin{lem}\label{lem:lattice trace class}
For every fixed $t\in\RR$,
\eq{\left[\Lambda(X),\ee^{\ii tH}\right]\in\mathcal{S}_1.}
Consequently,
\eq{\left[
\ee^{\ii tH}\Lambda(X)\ee^{-\ii tH},\Lambda(X)
\right]\in\mathcal{S}_1.}
\end{lem}

\begin{lem}[The lattice trace formula]\label{lem:lattice trace formula}
For every $\ve>0$ and $t\in\RR$,
\eql{&2\pi\ii\tr\left(
P_F\left[
\ee^{\ii tH}\Lambda_\ve(X)\ee^{-\ii tH},
\Lambda_\ve(X)
\right]P_F
\right)\notag\\
&\qquad=
-\frac{1}{\pi}
\int_{p\in\Sigma_{E_F}}\dif{p}
\int_{k\notin\Sigma_{E_F}}\dif{k}\,
\frac{\sin\left(t(f(p)-f(k))\right)}
{1-2\ee^{-\ve}\cos(p-k)+\ee^{-2\ve}}\,.
\label{eq:regularized lattice trace}}
Moreover,
\eql{&2\pi\ii\tr\left(
P_F\left[
\ee^{\ii tH}\Lambda(X)\ee^{-\ii tH},
\Lambda(X)
\right]P_F
\right)\notag\\
&\qquad=
-\frac{1}{\pi}
\int_{p\in\Sigma_{E_F}}\dif{p}
\int_{k\notin\Sigma_{E_F}}\dif{k}\,
\frac{\sin\left(t(f(p)-f(k))\right)}
{\left|\ee^{\ii p}-\ee^{\ii k}\right|^2}\,.
\label{eq:sharp lattice trace}}
Consequently, the sharp trace in \cref{eq:sharp lattice trace} is the
$\ve\to0^+$ limit of \cref{eq:regularized lattice trace}.
\end{lem}

\begin{proof}
With respect to the measure $\dif{k}/(2\pi)$, the momentum-space kernel of
the regularized switch is the convergent geometric series
\eq{\widehat{\Lambda_\ve}(p,k)
=
\sum_{n\geq0}\ee^{-\ve n}\ee^{-\ii n(p-k)}
=
\frac{1}{1-\ee^{-\ve-\ii(p-k)}}\,.}
Therefore,
\eq{\widehat{\Lambda_\ve}(p,k)
\widehat{\Lambda_\ve}(k,p)
=
\frac{1}
{1-2\ee^{-\ve}\cos(p-k)+\ee^{-2\ve}}\,.}
Since $\Lambda_\ve(X)$ is trace-class, the trace may be computed directly
from the kernels. This gives
\eq{&2\pi\ii\tr\left(
P_F\left[
\ee^{\ii tH}\Lambda_\ve(X)\ee^{-\ii tH},
\Lambda_\ve(X)
\right]P_F
\right)\\
&\qquad=
-\frac{1}{\pi}
\int_{p\in\Sigma_{E_F}}\dif{p}
\int_{k\in\TT}\dif{k}\,
\frac{\sin\left(t(f(p)-f(k))\right)}
{1-2\ee^{-\ve}\cos(p-k)+\ee^{-2\ve}}\,.}
The integral over $\Sigma_{E_F}\times\Sigma_{E_F}$ vanishes by antisymmetry
under $p\leftrightarrow k$, proving
\cref{eq:regularized lattice trace}.

We next identify its limit with the honest sharp trace. Under the Fourier
transform, the sharp switch is the Hardy projection with boundary kernel
$\sum_{n\geq0}\ee^{-\ii n(p-k)}$. Although this kernel is distributional,
the trace-class difference from \cref{lem:lattice trace class} has the
ordinary kernel
\eql{\left(
\ee^{\ii tH}\Lambda(X)\ee^{-\ii tH}-\Lambda(X)
\right)(p,k)
=
\frac{\ee^{\ii t(f(p)-f(k))}-1}
{1-\ee^{-\ii(p-k)}}\,.
\label{eq:lattice difference kernel}}
The apparent singularity at $p=k$ is removable.

Cyclicity of the trace, now justified by
\cref{lem:lattice trace class}, gives
\eq{&\tr\left(
P_F\left[
\ee^{\ii tH}\Lambda(X)\ee^{-\ii tH},\Lambda(X)
\right]P_F
\right)\\
&\qquad=
\tr\left(
\left(
\ee^{\ii tH}\Lambda(X)\ee^{-\ii tH}-\Lambda(X)
\right)
[\Lambda(X),P_F]
\right)\,.}
To evaluate the last trace, one may replace the second occurrence of
$\Lambda(X)$ by $\Lambda_\ve(X)$ and then let $\ve\to0^+$. This is legitimate
because $\Lambda_\ve(X)\to\Lambda(X)$ strongly, the commutators are uniformly
bounded, and the first factor is trace-class. Using
\cref{eq:lattice difference kernel} and then exchanging $p$ and $k$ in one
of the two cross terms yields
\eq{&\tr\left(
P_F\left[
\ee^{\ii tH}\Lambda(X)\ee^{-\ii tH},\Lambda(X)
\right]P_F
\right)\\
&\qquad=
\frac{\ii}{2\pi^2}
\int_{p\in\Sigma_{E_F}}\dif{p}
\int_{k\notin\Sigma_{E_F}}\dif{k}\,
\frac{\sin\left(t(f(p)-f(k))\right)}
{\left|1-\ee^{-\ii(p-k)}\right|^2}\,.}
Since
\eq{\left|1-\ee^{-\ii(p-k)}\right|^2
=\left|\ee^{\ii p}-\ee^{\ii k}\right|^2,}
this proves \cref{eq:sharp lattice trace}.

It remains to justify the limit in the regularized formula. Near the
diagonal,
\eq{1-2\ee^{-\ve}\cos(p-k)+\ee^{-2\ve}
=
(1-\ee^{-\ve})^2
+2\ee^{-\ve}(1-\cos(p-k))
\geq c\left(\ve^2+|p-k|^2\right).}
Furthermore,
\eq{\left|\sin\left(t(f(p)-f(k))\right)\right|
\leq C_t|p-k|\,.}
After the integral over
$\Sigma_{E_F}\times\Sigma_{E_F}$ has been removed, the remaining domain
meets the diagonal only at the finitely many endpoints of the Fermi arcs.
There the resulting bound $C_t/|p-k|$ is integrable across the corresponding
two-dimensional corner. Dominated convergence therefore gives
\eq{\lim_{\ve\to0^+}
\frac{1}
{1-2\ee^{-\ve}\cos(p-k)+\ee^{-2\ve}}
=
\frac{1}{|\ee^{\ii p}-\ee^{\ii k}|^2}}
inside the trace integral and proves the final assertion.
\end{proof}

\end{document}
